\documentclass[10pt]{article}
\usepackage[margin=0.8in]{geometry}
\usepackage{booktabs,array,url,hyperref}
\usepackage{microtype}
\hypersetup{hidelinks}
\title{Supplementary Computational Provenance\\Boundary-Localized Genomic Utility and Clinical-Tail Dominance}
\author{Ishanu Chattopadhyay}
\date{}
\begin{document}
\maketitle

\section*{Purpose}
The main manuscript describes analyses in statistical rather than implementation-specific terms. This supplement records the computational provenance required to regenerate the figures and tables without introducing implementation filenames into the scientific narrative.

\section*{Figure-generation policy}
All empirical manuscript figures are generated at compilation time with PGFPlots from tabular outputs in the repository analysis-results tree. No raster result figure is embedded in the manuscript. The figure source defines a single root path,
\begin{center}
\path{../zebra_comp/RESULTS},
\end{center}
so re-running the analyses and recompiling the manuscript updates plotted values without creating a separate manuscript-side data directory.

\section*{Primary source mapping}
\begin{tabular}{p{0.30\textwidth}p{0.64\textwidth}}
\toprule
Manuscript element & Analysis-output source under \path{zebra_comp/RESULTS} \\
\midrule
Global repeated-split discrimination & \path{032_REPEATED_SPLITS_03_TARGET_LOGIC/auc_distribution_summary.csv} and \path{paired_auc_delta_summary.csv} \\
Paired incremental-value figure & \path{033_INCREMENTAL_LOGISTIC_32_COHORT/paired_incremental_summary.csv} \\
ZeBRA--MUC5B stratification & \path{MUC5B_STRATIFIED_BY_FILD/MUC5B_ZEBRA_STRATIFIED_SUMMARY.csv} \\
Local information figure & \path{LOCAL_ZEBRA_MUC5B_INFORMATION_CURVES/FPR_DEFINED_LOCAL_INFORMATION_BANDS.csv} and \path{LOCAL_ZEBRA_MUC5B_INFORMATION_WINDOWS.csv} \\
Boundary rescue figure & \path{MUC5B_ZEBRA_RESCUE_RULE/REPEATED_SPLIT_MUC5B_RESCUE_SUMMARY.csv} \\
Stringent operating points & \path{032_REPEATED_SPLITS_03_TARGET_LOGIC/zedstat_operating_points_summary.csv} \\
Direct likelihood-ratio robustness figure & \path{EMPIRICAL_ZEBRA_MUC5B_LR_ROBUSTNESS/LR_ESTIMATOR_SENSITIVITY_GRID.csv}, \path{CLINICAL_TAIL_ROBUSTNESS.csv}, and \path{ROBUSTNESS_RANGE_SUMMARY.csv} \\
\bottomrule
\end{tabular}

\section*{Direct empirical likelihood-ratio analysis}
The score-level likelihood-ratio analysis is generated by
\begin{center}
\path{test_empirical_zebra_muc5b_likelihood_ratios.py}
\end{center}
and the estimator-sensitivity analysis by
\begin{center}
\path{test_empirical_zebra_muc5b_lr_robustness.py}.
\end{center}
Both use the same processed cohort inputs as the other analyses in \path{zebra_comp}. The primary direct-LR specification uses five spline knots and logistic regularization $C=1$. The prespecified robustness grid crosses spline knots $\{4,5,6\}$ with $C\in\{0.3,1,3\}$; it is a sensitivity analysis rather than a model-selection exercise.

\section*{Direct likelihood-ratio robustness table}
\begin{table}[!t]
\caption{Direct empirical likelihood-ratio summaries. The primary column uses the five-knot, $C=1$ estimator; the range column reports the minimum and maximum across the nine prespecified spline and regularization specifications.}
\label{tab:empirical-lr-summary}\centering\small
\begin{tabular}{lcc}
\toprule
Quantity & Primary value & Robustness range\\
\midrule
$b_{0.05}$ & 0.8699 & 0.8441--0.9083\\
$b_{0.01}$ & 0.9099 & 0.8635--0.9617\\
$b_{0.005}$ & 0.9241 & 0.8697--1.6100\\
Median $\log \Lambda_Z$ above 95th percentile & 1.5175 & 1.4383--1.6069\\
95th pct. residual genomic scale above 95th percentile & 0.9213 & 0.8155--5.6485\\
Fraction with $\log\Lambda_Z$ exceeding genomic scale above 95th percentile & 0.9045 & 0.8811--1.0000\\
Median $\log \Lambda_Z$ above 97.5th percentile & 2.7325 & 2.3756--2.8043\\
95th pct. residual genomic scale above 97.5th percentile & 0.7758 & 0.7601--8.4206\\
Fraction with $\log\Lambda_Z$ exceeding genomic scale above 97.5th percentile & 0.9938 & 0.8250--1.0000\\
Median $\log \Lambda_Z$ above 99th percentile & 3.8345 & 3.7850--3.9379\\
95th pct. residual genomic scale above 99th percentile & 0.7384 & 0.6849--0.7473\\
Fraction with $\log\Lambda_Z$ exceeding genomic scale above 99th percentile & 1.0000 & 1.0000--1.0000\\
\bottomrule
\end{tabular}
\end{table}

\section*{Exploratory local likelihood-ratio statistics}
The local FPR-band analyses are used descriptively in the main manuscript because multiple regions are examined within the same cohort. For reproducibility, the nominal one-degree-of-freedom likelihood-ratio-test p-values are recorded here but are not treated as multiplicity-adjusted confirmatory inference.
\begin{center}
\begin{tabular}{lrr}
\toprule
ZeBRA control-FPR band & LRT for adding \textit{MUC5B} after ZeBRA & Nominal $p$\\
\midrule
2--5\% & 4.44 & 0.0352\\
5--10\% & 16.55 & $4.74\times10^{-5}$\\
10--20\% & 4.77 & 0.0289\\
20--40\% & 13.69 & $2.15\times10^{-4}$\\
\bottomrule
\end{tabular}
\end{center}

\section*{Theory--data scope}
The empirical residual-genomic quantiles are fixed-score quantities for the observed ZeBRA representation and \textit{MUC5B}. They correspond to the finite-stage $b_{\delta,n}$ construction in Theorem~1(ii), not to the optional uniform-tightness condition A1-U over all clinical-history lengths. Likewise, the observed high-ZeBRA likelihood-ratio tail does not establish the mathematical unboundedness assumption A2.

\section*{Calibration summary}
The raw ZeBRA score is not an absolute calibrated FILD/FILA probability in this cohort. Across repeated held-out splits, the mean calibration slope is 0.180 for ZeBRA, 0.241 for the genomic model, and 1.140 for the combined model; corresponding expected calibration errors are 0.570, 0.115, and 0.121. The main manuscript therefore treats ZeBRA as an ordinal clinical-history score unless an explicit score-level likelihood-ratio transformation is estimated.

\section*{Supporting ROC-hull summary}
The descriptive ROC-hull analysis is retained as a supporting result rather than a primary predictive comparison. The empirical/interpolated AUCs are 0.8142 for ZeBRA and 0.8095 for the hybrid model; the convex hulls of the individual curves have AUCs 0.8219 and 0.8182, respectively; the pointwise maximum has AUC 0.8196; and the convex hull of the concatenated ZeBRA and hybrid ROC points has AUC 0.8260. These hull quantities are attainable geometric envelopes and are not held-out performance estimates from a separately fitted classifier.

\section*{Authoritative manuscript sources}
The canonical manuscript master is \path{zebra_genomics_boundary_theorem_empirical.tex}. The scientific narrative is split across \path{manuscript_body_intro.tex}, \path{manuscript_body_theory.tex}, \path{manuscript_body_interpretation.tex}, \path{manuscript_body_empirical_a.tex}, \path{manuscript_body_empirical_b.tex}, and \path{manuscript_body_end.tex}. Figure definitions are assembled by \path{empirical_figures.tex} from \path{empirical_figures_data.tex}, \path{empirical_figures_primary.tex}, and \path{empirical_figures_lr.tex}. These figure-source files are the authoritative mapping between manuscript plots and analysis-result tables.

\end{document}
